%% Research design section 3 -- the physical graph and the communication demand
%% expressed together.
%%
%% graphsearch/schema.py (normalisation: bytes/s throughout, directed edges,
%% shared resources as vertices), graphsearch/paths.py (allowed paths, cut
%% capacity, the boundary), graphsearch/demand.py (CommFlow).
%%
%% The unit trap of section 3 belongs here and is not a footnote: a link's
%% datasheet gbps, a collective's busbw and the simulator's link_bw answer
%% different questions, and heteropilot's A40 bridge measures 25.0 / 19.29 /
%% 8.8 GB/s for one wire depending on which is asked.
\section{Representing the cluster and its traffic}
\label{sec:graph}

\paragraph{One graph, normalised.}
The cluster becomes a directed graph whose vertices are accelerators, network
interfaces, CPU sockets and switches, and whose edges are links. Every capacity
is bytes per second, whatever unit the source file used; a schema that lets one
file mean gigabytes and another gigabits is a schema in which a factor of eight
survives review. Undirected links in the source become two directed edges, so a
path is a sequence of edges and not a sequence of ambiguous adjacencies.

\paragraph{Shared resources are vertices, not annotations.}
A resource that several links contend for---a host bridge, an uplink, a switch
port---enters the graph as its own vertex, with the links that cross it
incident on it. It carries a capacity and the share already reserved by traffic
outside this deployment. Making it a vertex rather than a label is what lets a
cut computation see it: the capacity between two device sets is then a maximum
flow in a graph where the contended element is a real node with a real
capacity, and a reservation is subtracted before any bound is computed rather
than after.

\paragraph{Demand.}
A candidate's traffic is a set of flows, each with its participants, its bytes
per event, its events per request, and whether a latency target charges for it.
The last field is load-bearing: ingress and egress cross an uplink and consume
its capacity, but no declared SLO charges for their time, so a model that
priced only what the SLO charges for would conclude that those flows are free
and that the uplink they saturate is idle.

\paragraph{A link does not have \emph{a} bandwidth.}
This is the trap the representation exists to avoid, and it is not a footnote.
One wire answers different questions with different numbers depending on what
crosses it. Measured on the eight-GPU node used in
Section~\ref{sec:eval}, one PCIe path between two accelerators sustains
\egfourPeerBridge~GB/s for a direct peer copy, \egfourBusbwTwo~GB/s of bus
bandwidth for a two-rank all-reduce, and \egfourBusbwFour~GB/s for the
four-rank all-reduce a tensor-parallel group of four actually runs.

% claim: C17
Those are three measurements of one path, and the last two differ by more than
a factor of two. A measurement is therefore keyed by what was run---the collective,
the number of ranks, the message-size band and how the measuring process was
bound to the node---and a figure recorded under one key is never substituted for
another. The number of ranks in particular is part of the key and not a note
beside it, because filing a two-rank figure where a four-rank one belongs is
precisely the substitution that keying exists to prevent.

% claim: C15
% claim: C16
\paragraph{Which resource is shared is settled by measurement.}
The representation can express that two links contend for something; it cannot
decide whether they do. That is a claim about hardware and it has to be
measured. Section~\ref{sec:eval} reports a case where the declaration was
written first and the measurement contradicted it: two device pairs believed to
share a host bridge were measured to share nothing, each concurrent copy
sustaining the full single-copy rate. The consequence for this section is
narrow and worth stating plainly---a shared-resource vertex is a hypothesis
about the machine, and it is wrong until something has put a probe on the wire.
